\documentclass[11pt]{article}
\usepackage[a4paper,margin=25mm]{geometry}
\usepackage{amsmath,amssymb,amsthm,mathtools,booktabs,array,microtype}
\usepackage[T1]{fontenc}
\usepackage{lmodern}
\usepackage[colorlinks=true,linkcolor=blue!40!black,urlcolor=blue!50!black,citecolor=blue!40!black]{hyperref}
\usepackage{xcolor}
\hypersetup{pdftitle={Three-scale extremal reconstruction of Dirichlet L-function zeros},pdfauthor={Research note}}
\newtheorem{theorem}{Theorem}[section]
\newtheorem{lemma}[theorem]{Lemma}
\newtheorem{corollary}[theorem]{Corollary}
\newtheorem{proposition}[theorem]{Proposition}
\theoremstyle{remark}\newtheorem{remark}[theorem]{Remark}
\newcommand{\N}{\mathbb N}\newcommand{\R}{\mathbb R}\newcommand{\C}{\mathbb C}
\newcommand{\Rea}{\operatorname{Re}}\newcommand{\ord}{\operatorname{ord}}
\newcommand{\limsupn}{\limsup_{n\to\infty}}
\newcommand{\bfa}{\boldsymbol\alpha}
\newcommand{\A}{\mathcal A}
\newcommand{\bb}{\boldsymbol b}
\title{Three-scale extremal reconstruction of Dirichlet $L$-function zeros\\[4pt]\large Sharp one-sided prime-sum limits on an integer lattice}
\author{Research note}
\date{26 November 2026}
\begin{document}
\maketitle
\begin{abstract}
Assuming the Dirichlet GRH theorem established in the supplied Lean development, we prove that a countable family of sharp upper limits of integer-coefficient filtered von Mangoldt sums determines the complete nontrivial zero multiset of every Dirichlet $L$-function, or of every conjugation-symmetric positive aggregate of such functions. All limits are taken on $x=6^n$, using finitely many neighboring $2$-$3$ smooth integers at each filter degree. The three dilation operators at scales $2,3,6$ simultaneously eliminate the main terms, preserve nonnegative zero weights, and separate all nonnegative zero ordinates. A compact moment problem then yields an explicit pointwise multiplicity formula, including central zeros and coincidences between distinct functions. No simplicity or linear-independence hypothesis is needed. The analytic inputs are classical; the contribution is the constant-free, sparse, one-sided reconstruction scheme and its universal two-scale obstruction.
\end{abstract}

\section{Source analysis and scope}
The supplied repository \cite{repo} is a rewrite of the Alp\"oge--Claude proof of RH and Dirichlet GRH. I retrieved the root README, the readable edition's theorem statements and analytic endgame files, and the compact edition's guide via the GitHub connector, at commit
\begin{center}\small\texttt{72fead102f663d5c7a407964a13dc0cf9d977207}.\end{center}
The repository reports Lean 4.35.0, Mathlib revision \href{https://github.com/leanprover-community/mathlib4/commit/56a005855ea17e0b0958f73ec801ae140b0a1679}{\texttt{56a00585}}, and independent checks using nanoda, con-leche, con-ron and Comparator. It reports only the ordinary Lean axioms \texttt{propext}, \texttt{Classical.choice} and \texttt{Quot.sound}. These are reported results, not a fresh replay in this investigation. Its root README explicitly distinguishes machine verification from human review. The supplied Wikipedia article \cite{wiki} likewise describes the November 2026 announcement and the prime-distribution consequences. We use the zero-line conclusion as the mathematical input, and not a claim about all automorphic $L$-functions.

\paragraph{Relevant architecture.} The compact guide separates three geometric/arithmetic bounds (short forms, enough forms, and long integral forms) assembled into a main inequality. In the readable edition, the analytic reduction uses a ramp-weighted prime-power remainder. In particular, \path{New/Reduction/Zeros/PrimeError.lean} studies, for an auxiliary quadratic character $\chi_d$,
\[
 R_{d,q}(x)=\sum_{\substack{m\le x\\m\equiv1\ (q)}}\Lambda(m)(1-\chi_d(m))
 (m^{-1}-x^{-1})-\frac{\log x}{\varphi(q)}-c_{d,q}.
\]
Its transform has residues proportional to the difference of multiplicities of zeros of the untwisted and twisted $L$-products, divided by $s(1-s)$. The file \path{New/Reduction/Zeros/Criterion.lean} uses suitable twists and a Landau argument to exclude zeros from a right half-plane using an eventual upper bound; see also Suzuki \cite{suzuki}. Finally \texttt{src/Main.lean} exports the nonvanishing and critical-line statements. The formal theorem \texttt{LFunction\_eq\_zero\_re} covers every positive modulus and character, primitive or imprimitive. The file \texttt{Challenge.lean} contains placeholder statements for the comparator, not a dependency of the proof according to the README.

Our approach runs in a complementary direction: once zeros are on the line, the residue denominator becomes the strictly positive number $1/4+\gamma^2$. Instead of subtracting twists (which can create signed spectral weights), we take nonnegative conjugation-symmetric character aggregates, and use dilation differences to remove the pole at $1$. We recover the spectral data from one-sided extrema, not from Fourier coefficients of the full prime-counting function.

\section{Statement of the reconstruction theorem}
Fix $q\ge1$. Let $\chi$ range over the Dirichlet characters modulo $q$, extended by zero at nonunits, and let $b_\chi\ge0$ satisfy $b_{\bar\chi}=b_\chi$. Set
\begin{align}
 w_{\bb}(m)&=\sum_{\chi\ (q)} b_\chi\chi(m)\in\R,\\
 P_{\bb}(x)&=\sum_{1\le m\le x}\Lambda(m)w_{\bb}(m)(x/m-1),\qquad x>0.
 \label{eq:prime}
\end{align}
Here $\Lambda(p^j)=\log p$ and $\Lambda(m)=0$ otherwise. The weights $w_{\bb}(m)$ need not be positive, although all the special choices in Section~\ref{sec:individual} are nonnegative. For the identity residue class choose $b_\chi=1/\varphi(q)$, so $w_{\bb}(m)=\mathbf 1_{m\equiv1\ (q)}$ by character orthogonality. For $q=1$, $P_{\bb}$ with $b_1=1$ is simply the ordinary ramp-weighted sum.

For $r\in\{2,3,6\}$ define the dilation difference
\begin{equation}\label{eq:operator}
 (\A_r f)(x)=(r+1)f(x)-f(rx)-r f(x/r),
 \qquad g_r(t)=r+1-2\sqrt r\cos(t\log r).
\end{equation}
The operators commute, and their coefficients are integers. Let
\begin{equation}\label{eq:moment}
 M_{ijk}(\bb)=\limsupn 6^{-n/2}
 (\A_2^{i+2}\A_3^j\A_6^k P_{\bb})(6^n),\qquad i,j,k\in\N_0.
\end{equation}
For $t\in\R$ let $m_\chi(t)$ be the order of vanishing of $L(s,\chi)$ at $s=1/2+it$, with zero order at nonzeros, and write $W_{\bb}(t)=\sum_\chi b_\chi m_\chi(t)$. Complex conjugation implies $W_{\bb}(-t)=W_{\bb}(t)$. Put $\eta(0)=1$ and $\eta(t)=2$ for $t>0$. Define the finite measure
\begin{equation}\label{eq:measure}
 \mu_{\bb}=\sum_{t\ge0}\frac{\eta(t)W_{\bb}(t)g_2(t)^2}{1/4+t^2}\,\delta_{g(t)},
 \qquad g(t)=(g_2(t),g_3(t),g_6(t)).
\end{equation}
The sum is over distinct nonnegative ordinates with nonzero weight; no simplicity assumption is made. Its support lies in the compact box
\[
 K=\prod_{r=2,3,6}[(\sqrt r-1)^2,(\sqrt r+1)^2].
\]

\begin{theorem}[Three-scale extremal reconstruction]\label{thm:main}
Under Dirichlet GRH, every limit superior in \eqref{eq:moment} is finite, and
\begin{equation}\label{eq:momentid}
 M_{ijk}(\bb)=\int_K u_2^i u_3^j u_6^k\,d\mu_{\bb}(u).
\end{equation}
The array $\{M_{ijk}(\bb)\}$ determines $W_{\bb}(t)$ at every real $t$. In particular, for every $t\ge0$ put
\begin{equation}\label{eq:localizer}
 Q_{N,t}(u)=\prod_{r=2,3,6}\left(1-\frac{(u_r-g_r(t))^2}{16r}\right)^N
 =\sum_{i,j,k}c_{ijk}(N,t)u_2^i u_3^j u_6^k.
\end{equation}
Then the following is an explicit inversion formula:
\begin{equation}\label{eq:recover}
 \boxed{\displaystyle
 W_{\bb}(t)=\frac{1/4+t^2}{\eta(t)g_2(t)^2}
 \lim_{N\to\infty}\sum_{i,j,k}c_{ijk}(N,t)M_{ijk}(\bb).}
\end{equation}
The inner sum is finite ($0\le i,j,k\le2N$). The unscaled sum in the limit is nonnegative and decreases to the desired atom mass. The result also holds if the moment sum is replaced by the single sharp upper limit for the filter $\A_2^2 Q_{N,t}(\A_2,\A_3,\A_6)$.
\end{theorem}

Thus one needs neither an unknown centering constant nor the values of the prime sum at all real arguments. For fixed $(i,j,k)$, each filter uses only a finite set of arguments $2^{n+a}3^{n+b}$, where $|a|\le i+2+k$ and $|b|\le j+k$. These are integers for sufficiently large $n$. The theorem concerns exact infinite limits, not a finite-data numerical algorithm.

\section{Proof: the filtered explicit formula}
We record the standard analytic inputs from classical Dirichlet $L$-function theory \cite{davenport,mv}: Euler products for $\Rea s>1$, meromorphic continuation (only the principal character has a simple pole at $1$), and the local zero-count and logarithmic-derivative estimates. For fixed $q$, the number of primitive nontrivial zeros in $T-1\le\Im s\le T+1$ is $O_q(\log(|T|+2))$, counted with multiplicity. On a fixed strip $1/4\le \sigma\le2$ at large $|T|$,
\begin{equation}\label{eq:localLD}
 \frac{L'}{L}(\sigma+iT,\chi)=
 \sum_{|\Im\rho-T|\le1}\frac{1}{\sigma+iT-\rho}
 +O_q(\log(|T|+2)),
\end{equation}
away from zeros, where the sum uses the primitive inducing function. For imprimitive characters the omitted Euler factors have zeros on $\Rea s=0$ and contribute bounded logarithmic derivatives on this strip. These facts follow from the Hadamard product, the gamma-factor estimates and the zero count; they are independent of GRH. For example, subtract the logarithmic derivative at $2+iT$ from that at $\sigma+iT$ in the Hadamard product. The nonlocal terms are bounded by a constant times $\sum_{|\gamma-T|>1}(\gamma-T)^{-2}=O_q(\log(|T|+2))$ by the local zero count, and the gamma terms are $O_q(\log(|T|+2))$. The logarithmic derivative at $2+iT$ is bounded by absolute convergence. This gives \eqref{eq:localLD}.

\begin{lemma}[Positive spectral expansion]\label{lem:formula}
Let $H(z_2,z_3,z_6)$ be any real polynomial. Under Dirichlet GRH, as $x\to\infty$,
\begin{multline}\label{eq:expansion}
 x^{-1/2}(\A_2^2 H(\A_2,\A_3,\A_6)P_{\bb})(x)\\
 =\sum_{t\ge0}\frac{\eta(t)W_{\bb}(t)g_2(t)^2}{1/4+t^2}
 H(g(t))\cos(t\log x)+O_{q,\bb,H}(x^{-1/4}).
\end{multline}
The sum is absolutely and uniformly convergent in $\log x$. If $H\ge0$ on $K$, all its cosine coefficients are nonnegative.
\end{lemma}
\begin{proof}
For $\Rea s>1$ the Euler products give
\[
 D_{\bb}(s):=\sum_{m\ge1}\frac{\Lambda(m)w_{\bb}(m)}{m^s}
 =-\sum_\chi b_\chi\frac{L'}{L}(s,\chi).
\]
Use the Mellin inversion identity, valid for $y>0$ and $c>1$,
\[
 (y-1)_+=\frac{1}{2\pi i}\int_{c-i\infty}^{c+i\infty}
 \frac{y^s}{s(s-1)}\,ds.
\]
Indeed its residues when closing to the left for $y>1$ are $y$ and $-1$; closing right for $y<1$ gives zero, and continuity covers $y=1$. Absolute convergence and the Dirichlet series imply
\begin{equation}\label{eq:perron}
 P_{\bb}(x)=\frac1{2\pi i}\int_{(2)}D_{\bb}(s)\frac{x^s}{s(s-1)}\,ds.
\end{equation}
Dilation differences act on $x^s$ by multiplication by
\[
 G_r(s)=r+1-r^s-r^{1-s},\qquad G_r(1)=G_r(0)=0.
\]
The filtered integral has additional factor $G_2(s)^2 H(G_2(s),G_3(s),G_6(s))$. It is bounded on the strip $1/4\le\Rea s\le2$ (the bound depends on $H$). The factor $G_2(s)^2$ cancels the possible double pole at $s=1$, including the constant-term ambiguity in the main term. Equivalently, $\A_2^2$ annihilates both $x$ and $x\log x$, since $\A_r x=0$ and $\A_r(x\log x)=(1-r)x\log r$.

Shift the filtered integral to $\Rea s=1/4$. By GRH and the nonvanishing for $\Rea s\ge1$ away from the principal pole, the only crossed poles are at $\rho=1/2+it$ with residue $-W_{\bb}(t)$ in $D_{\bb}$. On the new vertical line \eqref{eq:localLD} gives $O_q(\log(|T|+2))$ for $|T|\ge2$, since its distance from the line of zeros is $1/4$; on compact portions there are no poles. Choose horizontal heights tending to $+\infty$ and $-\infty$ at distance at least $c_q/\log(|T|+2)$ from every ordinate of a zero of every inducing character: in each unit interval, excluding intervals of twice that radius removes less than the entire interval by the local zero count, for sufficiently small $c_q>0$. On these heights \eqref{eq:localLD} gives $O_q(\log^2(|T|+2))$ uniformly across the strip. The horizontal integrals tend to zero for fixed $x$, and the new vertical integral is $O_{q,\bb,H}(x^{1/4})$ by absolute integrability against $(1+T^2)^{-1}$. The residue at $\rho$ is
\[
 \frac{W_{\bb}(t)}{\rho(1-\rho)} x^\rho
 G_2(\rho)^2 H(G_2(\rho),G_3(\rho),G_6(\rho)).
\]
Here $\rho(1-\rho)=1/4+t^2$, and $G_r(\rho)=g_r(t)>0$. The local zero count implies $\sum_t W_{\bb}(t)/(1/4+t^2)<\infty$, including any finite central order. All filter multipliers are bounded for fixed $H$. This proves absolute convergence and allows arbitrary order of summation. Dividing by $\sqrt x$ and pairing $t$ with $-t$ gives \eqref{eq:expansion}.
\end{proof}

\section{Proof: sharp one-sided bounds without independence}
\begin{lemma}[Simultaneous recurrence on a discrete grid]\label{lem:recurrence}
For real $\theta_1,\ldots,\theta_J$, every $\epsilon>0$ and every integer $B$ there is an integer $n>B$ with $|e^{in\theta_j}-1|<\epsilon$ for all $j$.
\end{lemma}
\begin{proof}
The sequence $(e^{im\theta_1},\ldots,e^{im\theta_J})$ for $m\ge1$ lies in a compact torus. Choose a convergent subsequence indexed by strictly increasing integers $m_a$. Fix one sufficiently late index $a$, then choose $b>a$ with $m_b-m_a>B$. Both tuples are arbitrarily close to their common limit, so the componentwise quotient, which is the tuple at $n=m_b-m_a$, is arbitrarily close to the identity.
\end{proof}
\begin{proposition}[All nonnegative filters have exact upper endpoints]\label{prop:endpoint}
If $H$ is a real polynomial nonnegative on $K$, then
\begin{equation}\label{eq:sharp}
 \limsupn 6^{-n/2}(\A_2^2 H(\A_2,\A_3,\A_6)P_{\bb})(6^n)
 =\int_K H(u)\,d\mu_{\bb}(u).
\end{equation}
The same value is the limsup as $x\to\infty$ over all real $x$. One may replace the geometric base $6$ by any fixed real $v>1$, although integer sampling of all the shifts is then not automatic.
\end{proposition}
\begin{proof}
By Lemma~\ref{lem:formula} the expression is $\sum_{t\ge0} a_t\cos(nt\log6)+o(1)$ with $a_t\ge0$ and $\sum_t a_t=A<\infty$. Thus the upper limit is at most $A$. Given $\epsilon>0$ select finitely many $t$ with omitted mass below $\epsilon$. Apply Lemma~\ref{lem:recurrence} with angles $t\log6$ to make all retained cosines greater than $1-\epsilon$ at arbitrarily large $n$. The omitted part is bounded below by minus its mass. Consequently the expression is at least $A-\epsilon A-2\epsilon+o(1)$ along arbitrarily large indices. Send $\epsilon$ to zero. The continuous limsup cannot exceed $A$ and cannot be smaller than its discrete sub-limsup. The proof with $\log v$ is identical.
\end{proof}
Choosing $H(u)=u_2^i u_3^j u_6^k$ proves \eqref{eq:momentid}. The exact endpoint is a special feature of nonnegative spectral weights. It would not follow from the same proof for arbitrary signed coefficients, such as those in an untwisted-minus-twisted difference.

\section{Proof: separating frequencies and recovering atoms}
\begin{lemma}[Three-scale separation]\label{lem:separate}
The map $g:[0,\infty)\to K$ in \eqref{eq:measure} is injective.
\end{lemma}
\begin{proof}
Write $a=\log2$, $b=\log3$. Their ratio is irrational: an integer relation $pa=qb$ with positive integers would imply $2^p=3^q$. If $g(t)=g(t')$, then $\cos(at)=\cos(at')$, $\cos(bt)=\cos(bt')$ and $\cos((a+b)t)=\cos((a+b)t')$. Put $A=at$, $B=bt$, and $A'=at'$, $B'=bt'$. The first two equalities say $A'\equiv\varepsilon A$ and $B'\equiv\delta B\pmod{2\pi}$ for signs $\varepsilon,\delta\in\{1,-1\}$. If the signs agree, we have both congruences with the same sign. If they disagree, the third equality gives $\cos(A+B)=\cos(A-B)$, hence $\sin A\sin B=0$. For an angle with sine zero, its positive and negative signs agree modulo $2\pi$, so again both signs can be chosen equal, say $\varepsilon$. Therefore $a(t'-\varepsilon t)$ and $b(t'-\varepsilon t)$ are both integer multiples of $2\pi$. Irrationality forces $t'=\varepsilon t$. Nonnegativity gives $t'=t$.
\end{proof}
\begin{proof}[Completion of Theorem~\ref{thm:main}]
The measure in \eqref{eq:measure} is finite by the local zero count and $g_2$ bounded on the real axis. Every two coordinates $u_r,g_r(t)$ in the corresponding interval of $K$ differ by at most $4\sqrt r$. Thus each factor of \eqref{eq:localizer} is in $[0,1]$. At $u=g(t)$ all factors are $1$; at any other point some factor is strictly less than $1$. So $Q_{N,t}$ decreases pointwise to the indicator of the singleton $\{g(t)\}$. Dominated convergence and \eqref{eq:momentid} give
\[
 \lim_{N\to\infty}\sum_{i,j,k}c_{ijk}(N,t)M_{ijk}(\bb)
 =\mu_{\bb}(\{g(t)\})=\frac{\eta(t) W_{\bb}(t)g_2(t)^2}{1/4+t^2},
\]
where the last step uses Lemma~\ref{lem:separate}. Since $g_2(t)\ge(\sqrt2-1)^2>0$, division is legitimate even at $t=0$. Proposition~\ref{prop:endpoint} applied to $H=Q_{N,t}$ proves the alternative one-filter formula. This completes the proof.\qedhere
\end{proof}

\begin{corollary}[Determination by positive moments]
For fixed $(q,\bb)$ the moment array determines the finite Borel measure $\mu_{\bb}$ on $K$ uniquely. Every polynomial nonnegative on $K$ supplies an exact one-sided bound via \eqref{eq:sharp}. In particular, the moment matrices indexed by any finite set of monomials are positive semidefinite:
\[
 \sum_{\bfa,\boldsymbol\beta}v_{\bfa}v_{\boldsymbol\beta}
 M_{\bfa+\boldsymbol\beta}\ge0.
\]
\end{corollary}
\begin{proof}
Real polynomials are uniformly dense in $C(K)$ by Stone--Weierstrass. Equality of all moments of two finite measures therefore implies equality of integrals of all continuous functions and hence equality of the measures. The quadratic form is the integral of a polynomial square. The endpoint identity is Proposition~\ref{prop:endpoint}.
\end{proof}

\section{Individual \texorpdfstring{$L$}{L}-functions and multiplicities}\label{sec:individual}
The identity-class sum recovers the total orders $\varphi(q)^{-1}\sum_\chi m_\chi(t)$. To recover individual real characters, or complex characters up to conjugation, one can choose nonnegative arithmetic weights as follows. The principal character is recovered with $w_0=\chi_0$. For any nonprincipal $\chi$, take $w_\chi=\chi_0+\Rea\chi$. This is a nonnegative real weight: it vanishes on nonunits and is in $[0,2]$ on units. The character coefficients are nonnegative and symmetric: $b_{\chi_0}=1$ and $b_\chi=1$ for real $\chi$, or $b_\chi=b_{\bar\chi}=1/2$ for nonreal $\chi$.
Subtract the recovered $W_0(t)$ from $W_\chi(t)$ (not the individual limsup values of an arbitrary signed series). The resulting order is $m_\chi(t)$ for real characters and $(m_\chi(t)+m_{\bar\chi}(t))/2$ otherwise. The latter is the symmetrized spectrum of one conjugate pair. We now remove the remaining ambiguity using the complex character itself and a phase-sensitive filter.

\begin{corollary}[Primitive-conductor reconstruction]
Let $W_d^{\mathrm{id}}(t)$ denote the recovered weight from the identity-class sum modulo $d$, with every character coefficient $1/\varphi(d)$. For every positive $q$, the aggregate order at $1/2+it$ of all primitive characters of conductor exactly $q$ is
\[
 \sum_{\substack{\chi\ (q)\\\chi\ \mathrm{primitive}}} m_\chi(t)
 =\sum_{d\mid q}\mu(q/d)\varphi(d)W_d^{\mathrm{id}}(t),
\]
where $\mu$ is the ordinary M\"obius function. Only unweighted identity-class prime sums for the divisors of $q$ are required.
\end{corollary}
\begin{proof}
Each character modulo $q$ has a unique inducing primitive character of conductor dividing $q$, and every such primitive character induces exactly one character modulo $q$. The additional Euler factors have no zeros on the critical line, so induction preserves the multiplicities there. Thus $\varphi(q)W_q^{\mathrm{id}}(t)$ is the divisor sum of the primitive-conductor orders, and M\"obius inversion gives the result.
\end{proof}

\paragraph{Oriented extension.} For a nonreal character, use the complex-valued prime sum $P^{\chi}(x)=\sum_{m\le x}\Lambda(m)\chi(m)(x/m-1)$, which is obtained from the residue-class prime sums. The same contour proof, without pairing, gives
\[
 x^{-1/2}(\A_2^2 H(\A_2,\A_3,\A_6)P^\chi)(x)
 =\sum_{t\in\R}\frac{m_\chi(t)g_2(t)^2}{1/4+t^2}H(g(t))e^{it\log x}
 +O(x^{-1/4}).
\]
Here the notation $g(t)$ is extended evenly to all $t$. To resolve signs, fix a real target $t_0$ and choose $r\in\{2,3\}$ with $\sin(t_0\log r)\ne0$ if $t_0\ne0$; irrationality ensures at least one. Let $T_rf(x)=f(rx)$ and define the real-valued observable
\begin{align*}
 F_{N,t_0}(x)={}&\Rea\Big[x^{-1/2}\A_2^2 Q_{N,t_0}(\A_2,\A_3,\A_6)\,\\[-5pt]
 &\hspace{17mm}\cdot (2I-r^{-1/2}e^{-it_0\log r}T_r
       -r^{1/2}e^{it_0\log r}T_{1/r})P^\chi(x)\Big].
\end{align*}
The polynomial $Q_{N,t_0}$ means $Q_{N,|t_0|}$, and the last filter has multiplier $h(t)=2-2\cos((t-t_0)\log r)\ge0$ on $x^{1/2+it}$. By the same recurrence argument its sharp upper limit is the sum of the nonnegative spectral coefficients, even without symmetry of $m_\chi(t)$. As $N\to\infty$ this sum decreases to
\[
 \frac{g_2(t_0)^2}{1/4+t_0^2}\,m_\chi(-t_0)\,[2-2\cos(2t_0\log r)],
\]
with the coefficient in brackets strictly positive. Thus changing $t_0$ to $-t_0$ recovers $m_\chi(t_0)$ itself. At $t_0=0$ simply omit the last filter; the localizer alone gives $4g_2(0)^2m_\chi(0)$. This proves the stronger statement that the oriented zero multiset of every individual $L(s,\chi)$ can be reconstructed from real sharp upper limits of complex-character prime sums. In this extension the final filters may have complex coefficients before taking the real part, but the core dilation differences remain integer-coefficient operators. The conjugation-symmetric theorem itself needs only real weighted sums and integer-coefficient differences.

\section{Why use three scales? A sharp obstruction}
The two individual cosine coordinates need not be injective even though $\log2/\log3$ is irrational. This is not merely a possible numerical issue.
\begin{proposition}[Universal two-scale obstruction]\label{prop:alias}
For $a=\log2$, $b=\log3$, let
\[
 t=\pi(2/b-1/a),\qquad t'=\pi(2/b+1/a).
\]
Then $0<t<t'$, $g_2(t)=g_2(t')$, $g_3(t)=g_3(t')$, but $g_6(t)\ne g_6(t')$. In the class of arbitrary positive, summable cosine spectra, two-scale moment endpoints (with the fixed $g_2^2$ filter) cannot determine the spectrum.
\end{proposition}
\begin{proof}
The inequality $2/b>1/a$ is equivalent to $4>3$. Also $a(t'-t)=2\pi$ and $b(t'+t)=4\pi$, proving the two cosine equalities. The third equality is impossible by Lemma~\ref{lem:separate}. Place the same positive mass at either of the two distinct frequencies in the already filtered measure. Every polynomial in the first two coordinates has identical integrals and thus identical endpoints; the third coordinate separates these measures. Equivalently, before the fixed positive filter the two masses can be rescaled to compensate for its multiplier. This is an obstruction for general spectra; it does not assert that either of the exhibited frequencies is an actual $L$-zero or that two scales cannot work for the particular arithmetic spectra under additional hypotheses.
\end{proof}

\section{Relation to existing mathematics and further questions}
The ingredients (ramp Perron formula, character orthogonality, compact recurrence, positive moments) are classical. The particular conclusion is a reconstruction theorem from \emph{only sharp upper endpoints} of sparse, constant-free dilation filters, not a new explicit formula for $L$-functions. The fixed squared filter is crucial: it cancels both leading terms without killing any critical-line frequency. The dependent scale $6=2\cdot3$ resolves the sign ambiguity of the two independent scale angles. This allows a direct, monotone atom extraction using nonnegative polynomials even if ordinates coincide, are rationally dependent, or are zero.

A targeted literature search on one-sided smoothed prime sums, variants of Chebyshev's conjecture, and recovery of zero spectra from extremal or moment data found related work on sign criteria (notably Suzuki \cite{suzuki}) and the classical explicit-formula framework \cite{davenport,mv,wiki}. The specific three-scale sparse endpoint scheme above was not located in that search. Priority remains unverified; no claim is made that the component techniques are new, or that this resolves a previously named open conjecture.

\paragraph{Further research questions.} (1) Can the fixed three-scale construction be replaced by two different nonnegative filters, with a universal injective spectral map and the same sparsity? (2) Can effective recurrence estimates and a priori zero-density bounds give a useful finite-sample reconstruction of the first few zeros? High-degree filters amplify numerical cancellation, and the present proof gives no useful recurrence rate for infinitely many phases. (3) Can analogous positive-residue reconstruction apply to other $L$-functions known to satisfy their own RH? The Dirichlet GRH theorem does not imply general automorphic RH. (4) Can the finite algebraic filter identities and recurrence lemma be formalized as generic Lean lemmas, and connected to a formal Perron shift and a compact moment argument? The present note is a human proof using classical analytic inputs, not a new Lean-checked theorem.

\section{Reproducibility and limitations}
The accompanying Python program expands the integer filters using exact arithmetic, checks cancellation of the leading terms, checks the two-scale alias and the third-scale separation numerically, and evaluates selected filtered prime-power sums for $q=1$ at finite powers of $6$. These computations are illustrations and tests, not a proof of any infinite-limit identity. The mathematical proofs are in Sections 3--7. All constants in the error estimates may depend on the fixed modulus, character weights, and polynomial; no uniformity in the filter degree or modulus is asserted. The inversion takes the sharp upper limit first, then the localization limit. Interchanging them is not justified.

For reproducibility, the finite experiment for $q=1$ produced the following values. Each column is the filtered sample inside the limsup in \eqref{eq:moment}; no column is an approximation with a certified error to the eventual endpoint.
\begin{center}
\small
\begin{tabular}{rrrrrr}
\toprule
$n$ & $6^n$ & $(0,0,0)$ & $(1,0,0)$ & $(0,1,0)$ & $(0,0,1)$\\
\midrule
3 & 216 & 0.450620 & 2.383286 & 2.875939 & 1.965527\\
4 & 1296 & 0.163647 & 0.999599 & 1.589280 & $-0.579067$\\
5 & 7776 & 0.253443 & 1.377211 & 1.644412 & 0.682193\\
6 & 46656 & 0.282123 & 1.435948 & 1.744829 & 1.036639\\
\bottomrule
\end{tabular}
\end{center}
The program also verifies exactly the cancellation identities for $64$ different filters and verifies numerically the multiplier and alias examples. The full output is in \texttt{verification.txt}.

\begin{thebibliography}{9}
\bibitem{repo} L. Sernau et al., \emph{GRH for Dirichlet L-functions, in Lean 4}, 25 November 2026, revision \href{https://github.com/lsernau/grh-lean/tree/72fead102f663d5c7a407964a13dc0cf9d977207}{\texttt{72fead1}}. See \href{https://github.com/lsernau/grh-lean/blob/72fead102f663d5c7a407964a13dc0cf9d977207/grh/src/Statements.lean}{Statements.lean}, \href{https://github.com/lsernau/grh-lean/blob/72fead102f663d5c7a407964a13dc0cf9d977207/grh/src/New/Reduction/Zeros/PrimeError.lean}{PrimeError.lean}, and \href{https://github.com/lsernau/grh-lean/blob/72fead102f663d5c7a407964a13dc0cf9d977207/grh/src/New/Reduction/Zeros/Criterion.lean}{Criterion.lean}.
\bibitem{wiki} Wikipedia, \emph{Riemann hypothesis}, accessed 26 November 2026. \url{https://en.wikipedia.org/wiki/Riemann_hypothesis}.
\bibitem{suzuki} M. Suzuki, On variants of Chebyshev's conjecture, \emph{Ramanujan J.} \textbf{68} (2025), article 95. \href{https://doi.org/10.1007/s11139-025-01238-9}{doi:10.1007/s11139-025-01238-9}.
\bibitem{davenport} H. Davenport, \emph{Multiplicative Number Theory}, 3rd ed., revised by H. L. Montgomery, Springer, 2000. Classical Dirichlet $L$-functions, zeros, logarithmic derivatives, and explicit formulas.
\bibitem{mv} H. L. Montgomery and R. C. Vaughan, \emph{Multiplicative Number Theory I: Classical Theory}, Cambridge University Press, 2007. In particular the explicit formula and oscillation/Landau methods.
\end{thebibliography}
\end{document}
